\documentclass{report}
\usepackage{amsfonts, amsmath}
\newcommand\R{{\mathbb R}}
\newcommand\Q{{\mathbb Q}}
\newcommand\N{{\mathbb N}}
\newcommand{\vf}{\varphi}
\pagestyle{empty}
\begin{document}
\hfuzz 12pt
\large
\centerline{\Huge \bf Analisi Matematica I}
\vskip.2cm
\centerline{\Huge Correzione Quarto Appello}

\vskip1cm
\begin{enumerate}
\item
$\log_3 9=2$, $\log_25\geq \log_2 4=2$.
\item 
Si applichi l'algoritmo euclideo. Poiche l'ultimo divisore \`e uno i
due numeri sono primi tra loro.
\item $f(x)\geq 1/2\, x-1/4\, x\geq 0$, $f(x)\leq 3/4$, la stima sulla
derivata \`e banale. Ora sia $x_n=f^n(0.5)$, poiche $f(0)=0$ si ha (Lagrange)
\[
|x_{n+1}|=|f(x_n)-f(0)|\leq \frac 34 |x_n|.
\]
Iterando tale disuguaglianza si ottiene $x_n\leq (\frac 34)^n$.

\item Basta dividere nel caso $x\leq 1$ per cui $f(x)=\ln(1)=0$ e il
caso $x>1$ per cui $f(x)=\ln(2x-1)$. Il resto \`e la solita roba.
\item Usando Lagrange ripetutamente si ha
\[
\ln\left(\sqrt{\frac{n^3+n}{n^3+1}}\right)=\frac 1{n^2}+{\mathcal O}(\frac
1{n^3}).
\]
Da cui segue la convergenza della serie.
\item {\sc Sol. informale}: Ovviamente si tratta dei triangoli
isosceli. Visto che l'area \`e area per altezza diviso due, il massimo
si ottiene per l'altezza massima (cio\`e uguale ad uno).\newline
{\sc Sol. formale}: Sia $x$ l'altezza del triangolo, chiaramente $x\in
[0,1]$. Allora l'area \`e data da $\frac 12 2 x=x$. Chiaramente ilk
massimo si ha per $x=1$.
\end{enumerate}
\end{document}
